\chapter{Recomendações para Trabalhos Futuros}

A arquitetura desenvolvida permite diferentes extensões, que podem ampliar o conjunto de ensaios e as possibilidades de pesquisa. Como trabalhos futuros, recomendam-se:

\begin{itemize}
  \item expandir a reprodução e a aquisição para três tensões e três correntes, com pelo menos seis canais analógicos sincronizados, permitindo ensaios trifásicos, faltas assimétricas e cálculo de componentes simétricas;;
  \item avaliar cenários com eventos múltiplos, como faltas simultâneas ou combinações de falta, variação de tensão e alteração de carregamento, de modo a submeter a ESP32 a ensaios de estresse computacional e verificar sua capacidade de processamento, temporização, uso de memória e robustez dos algoritmos embarcados;;
  \item ampliar o processador COMTRADE para outros formatos e diferentes ordenações de canais;
  \item tornar o mapeamento dos GPIOs integralmente configurável pelo JSON;
  \item avaliar conversores analógico-digitais externos e taxas de amostragem superiores em estudos de transitórios e harmônicos;
  \item incorporar curvas inversas normalizadas IEC e IEEE à função 51;
  \item implementar comparadores direcionais angulares adicionais para a função 67;
  \item desenvolver métodos de localização de faltas que considerem resistência de falta e condições do sistema;
  \item acrescentar CRC e identificação dos arquivos ao protocolo de transferência;
  \item avaliar a prototipagem dos blocos de geração, aquisição e processamento em FPGA, ou em arquiteturas híbridas com microcontrolador e FPGA, permitindo maior paralelismo, sincronismo entre canais e capacidade de processamento para uma versão mais robusta da bancada; e
  \item desenvolver uma placa dedicada com condicionamento analógico, isolamento e conectores próprios para a bancada.
\end{itemize}

Essas recomendações dão continuidade às limitações identificadas na implementação e aos resultados obtidos nos ensaios, orientando a evolução do protótipo para uma bancada didática mais abrangente, reconfigurável e metrologicamente caracterizada.
